% ===== §15 Sharing: The Birth of the Relational =====
\section{The Birth of the Relational}
\label{p22:sec:sharing}

\epigraph{\zh{独乐乐，与人乐乐，孰乐？}}{\textit{Mencius}}

\noindent
Sharing is the hinge of the whole model, the step at which beauty passes from an experience within a subject into an event between subjects. The two steps before it, perception and reception, could in principle be completed by a single person alone, and a tradition of aesthetic education has in fact stopped there, cultivating a refined private taste and calling the cultivation complete. If the model stopped where they stop, it would be one more such tradition, an education of the solitary connoisseur. What forbids the stopping is the concept the model serves. If beauty is produced in the coupling, then a beauty perceived and received but never carried into the between is a beauty that has not yet become what relational aesthetics says beauty is, and the carrying-across is therefore not an optional generosity laid over a beauty already complete but the act in which beauty is completed, in which the relational is born. This section asks what that act is, and it finds, as the two before it found, that what looks like a simple thing, telling another what one has seen, is the meeting-point of several traditions that disagree about what sharing does and whether, at its limit, it preserves or dissolves the two who share.

\subsection{The semantic pivot: sharing generates rather than transmits}
\label{p22:sub:sharing-pivot}

Everything in this section turns on a distinction that must be drawn at the outset, because the wrong reading of sharing quietly undoes the concept the paper has built. On one reading, sharing is transmission: I have a beauty, complete in me, and I convey it to you, so that afterward we both possess the same thing, as one might copy a file. On this reading beauty is presupposed as prior, private, and portable, a stock I hold and then duplicate. On the other reading, sharing is generation: the beauty is not conveyed from me to you but produced, then and there, in the act of joint attention, so that before the sharing there was a perception in me and an appropriateness in the field, but the beauty as a coupling event did not yet exist and comes to be only in the between. The first reading keeps beauty in the subject and adds a social act afterward; it is the second party to the old quarrel with a postal service attached. The second reading is the one relational aesthetics requires, on which beauty is born in the sharing and not merely delivered by it.

\claim{Sharing generates beauty rather than transmitting it. The act of joint attention does not convey a beauty already complete in one subject to another; it is the event in which beauty, as a coupling between subjects, first comes to be. What precedes the sharing is a perception within a subject and an appropriateness in the field; the beauty proper, the relational event, is produced in the sharing and has no prior existence to be transmitted.}

The distinction is not verbal. The transmission reading makes the third step dispensable in principle, a convenience for spreading a good that could have remained whole in one head; the generation reading makes it constitutive, the step without which beauty in the relevant sense never occurs at all. And the generation reading sets the section's central problem, because if beauty is generated in joint attention then we must say what joint attention is and how it differs from two people merely reporting the same pleasure side by side, the side-by-side arrangement that the concept section already warned is not coupling but its counterfeit. The traditions gathered below are gathered to answer exactly this, what makes the between an event and not a coincidence of two privacies.

\begin{casebox}[The light on the wall, seen alone and seen together]
In the late afternoon a particular light falls on the wall of the room the two share, a contingency of weather and season that will not fall in quite this way again. One of them sees it and is moved, registers its passing appropriateness, receives it. Here the cycle could stop, and often does: the one who saw keeps it, a private beauty appreciated alone, and the other never knows the light was there. Or the one who saw says, simply and at once, look, and the other looks, and for a moment the two hold the light in a focus each knows the other shares. The claim of this section is that something is produced in the second case that does not exist in the first, and that it is not merely the first beauty now had by two. Before the saying, there was one person's perception of a light. After the looking-together, there is a beauty that belongs to neither alone, a thing the two now hold between them, which will enter the memory of the room and the relation in a way the privately appreciated light would not. Whether this is real, and what it is if it is real, is the question the traditions below are summoned to answer; the case is only to fix what is at stake, the difference between a light seen alone and the same light made, by the saying, into something shared.
\end{casebox}

\subsection{Joint attention and the structure of the between}
\label{p22:sub:sharing-joint-attention}

The first account comes from the study of how two consciousnesses come to share a single attended world. Joint attention is not two people happening to look at the same thing; it is the recursive structure in which each attends to the thing and to the other's attending, so that the object is held in a common focus that each knows to be common. The developmental literature treats this triadic structure, the two subjects and the shared object bound by mutual awareness of the sharing, as the achievement on which the distinctively human capacity for shared intentionality is built~\parencite{tomasello_origins}; the phenomenological analyses of pairing and of the relation in which two are addressed as we describe the same structure from within, the lived sense in which the object is not mine-then-also-yours but ours from the start. To share a beauty, on this account, is to bring the field into a common focus of this recursive kind, in which what is attended is constituted as common in the very act of attending.

This gives the section the structural answer it needs to the problem of coincidence. Two people reporting the same pleasure side by side have not entered joint attention; each has a private focus that happens to fall on the same object, and the recursive layer, each attending to the other's attending and knowing the focus to be shared, is absent. Joint attention is precisely the presence of that layer, and it is the layer that makes the between an event rather than a coincidence. The account also reaches forward to the implementation section, since joint attention, unlike a private judgement, is the kind of thing that should leave a signature in the coordination of two nervous systems, and the synchrony that hyperscanning studies look for is at least a candidate trace of the recursive common focus this account describes. The account is strong on structure. What it does not by itself supply is the dimension of weight, the sense in which sharing is not only a coordination of attention but a matter of consequence between the two, and for that a second tradition is needed.

\subsection{Witnessing: the ethical weight of the shared}
\label{p22:sub:sharing-witnessing}

The second account has run quietly through this series from the beginning, in the reading of poetry as a showing that witnesses and in the practice of witnessing that the series proposed for its own community. Its claim is that sharing, at its depth, is mutual witnessing: you witnessed that this beauty, this contingency, this passing appropriateness, occurred in me, and your witnessing gives it a reality it could not have had alone. The point is not that two observations are more accurate than one. It is that being witnessed changes the standing of what is witnessed, lifts it from a private occurrence that one bears alone into a shared reality that two now hold. The literature on the witnessing of suffering shows the structure at its starkest: an experience that is not witnessed remains, for the one who underwent it, strangely unreal, hard to integrate, and the arrival of a witness is what allows it to become an experience that can be borne and woven in. Carried over to beauty, the structure says that a beauty witnessed becomes a common reality of the relation in a way a beauty unwitnessed does not, and the witnessing carries an ethical weight, the weight of not being left alone with what one has undergone.

This raises the account from coordination to consequence, and it supplies what joint attention alone could not. Joint attention says how the common focus is structured; witnessing says why it matters, what is at stake between the two in the sharing, namely whether each is left alone with experience or accompanied in it. The two are not rivals. Joint attention is the structure within which witnessing occurs, and witnessing is the ethical dimension of that structure, the reason the recursive common focus is not a cognitive curiosity but a thing that bears on how two people stand to one another. Together they begin to say what the between is. What they have not yet said is how the between moves, how the shared beauty does not merely sit, jointly attended and mutually witnessed, but grows in the exchange, and for that a third tradition is required.

\subsection{Dialogue and resonance: the movement of the shared}
\label{p22:sub:sharing-dialogue}

The third account insists that meaning is not held but made, and made only in the back-and-forth of address and answer. On the dialogical view, meaning lives in responsivity; an utterance that meets no answer is not yet fully meaningful~\parencite{bakhtin_dialogic}, and a meaning monologically held, never answered, is a meaning dying. Sharing, under this account, is the answering exchange in which the shared beauty does not stay fixed but increases, each response altering what it responds to, so that what passes between the two is not conserved but multiplied. To this the sociology of resonance adds a contrast that the political economy section will need: resonance names a relation in which subject and world answer and alter one another, mutually moved, and it is opposed to reification~\parencite{rosa_resonance}, the relation of unanswered possession in which the world is held mute and used. Sharing that resonates is sharing in which the two are mutually moved and changed; sharing that reifies is the holding of a beauty as a mute possession, attended perhaps, even witnessed, but not allowed to answer back and move.

The value of this account is that it supplies the section's dynamism and so its bridge to reproduction. Joint attention and witnessing, taken alone, might leave the shared beauty static, a common object jointly held and mutually confirmed but unchanging. The dialogical and resonant account makes the sharing itself a movement in which the beauty increases through answering, and an increase through answering is already the beginning of the reproduction the next section studies, the thickening of the shared in the exchange. Its commitment is that meaning and beauty are dialogically generated and non-conserved, born in the answer and multiplied in the back-and-forth, and this commitment will both feed the account of reproduction and sharpen the conflict the section must now face, the conflict over whether the union that sharing seeks preserves the two or dissolves them.

\subsection{A sociological voice: the effervescence of the gathered}
\label{p22:sub:sharing-effervescence}

One further account belongs here, drawn from the sociology of the gathered group, because it describes the production of a shared reality in the assembly with a precision that scales down instructively to the couple. The account observes that when people gather and attend together, something is generated that none brought and none holds alone, a heightened common reality, an effervescence~\parencite{durkheim_elementary}, in which the gathered feel themselves lifted into a shared life larger than their separate ones, and that the symbols and meanings a community lives by are forged in such moments and carried away from them as the residue of the gathering. The mechanism is the mutual amplification of attention and feeling in co-presence: each sees the others moved, is moved further by seeing it, and the shared movement compounds into a reality that is genuinely collective, located in the assembly and not in any member. This is sharing at the scale of the crowd, and its lesson for the couple is that the production of a common reality in co-present joint attention is not a peculiarity of the intimate dyad but a general feature of gathered attention, observed and analysed at the scale of the festival and the rite, of which the couple's shared moment is the smallest instance. The light on the wall, shared, produces between two what the festival produces among many, a common reality forged in co-present mutual attention and carried away as the residue that thickens the relation, and the sociological account, by exhibiting the mechanism at the large scale, confirms that the small-scale version is no sentimental invention but the dyadic case of a documented social process.

The account also contributes a caution the section will need. The effervescence of the gathered is morally indifferent, as capable of forging a lynch mob's terrible solidarity as a festival's joy, and this is the orthogonality of refinement and justice returning at the scale of sharing: the production of a powerful common reality in joint attention is no guarantee of the justice of what is produced. A couple, like a crowd, can forge in their shared attention a common reality that is generative and just, or one that is generative and cruel, accumulating beautifully upon the exclusion or consumption of a third. The sociological voice, by showing effervescence to be indifferent to the justice of its product, reinforces the paper's insistence that the cultivation of sharing, like every face of the sensibility, must carry the separate demand of justice, since the mere production of a vivid common reality, however real and however moving, settles nothing about whether what the sharing builds is good.

\subsection{Music and the common: an Eastern principal voice}
\label{p22:sub:sharing-music}

The principal Eastern voice on sharing is the oldest technology a civilisation devised for turning private feeling into common accord, and it belongs to the ethical-political heart of the tradition this series draws on. Music, the Record of Music says, is what is held in common, that which makes for union~\parencite{yueji}, \zh{乐者为同}; it is the harmony of heaven and earth, and its work is to bring hearts into a shared resonance that is concord without uniformity, harmonising and yet not merging. Mencius asks whether the joy of music alone can equal the joy of music shared with others, and the question answers itself, that beauty unshared falls short of beauty that has entered the common. This voice states, more directly than any of the Western accounts, the thesis the whole section serves, that beauty must enter the relational to complete itself, and it states it as an ethical and political truth and not merely an aesthetic one, since the rite-and-music tradition understood the sharing of beauty as the very means by which a common life is tuned and a people's disposition formed.

The voice carries two further gifts. The first is the figure of concord-without-uniformity, harmony that does not flatten the many into one, which is precisely the resource the section needs for the conflict ahead, since it names a union that preserves difference rather than dissolving it. The second is the seamlessness with which it joins the aesthetic to the ethical and the political, so that the sharing of beauty is shown to be, from the start and without strain, the tuning of a shared life. Where the Western accounts arrive at the ethical weight of sharing as a result, this tradition began there, having never divided the beautiful from the good, and it thereby confirms from the Eastern side the descent the whole paper traces, the disclosure of the ethical beneath the aesthetic, here met not as a discovery but as an inheritance.

The contrast the tradition draws between music and rite is worth drawing out, because it states the structure of sharing with a precision the Western accounts approach only slowly. Rite, in the pairing, is what differentiates, what marks the distinctions of role and station and keeps each in its place; music is what unites, what brings the differentiated into a common accord without erasing the distinctions that rite has marked. The two are held together deliberately, and the holding is the whole teaching: union without differentiation would be a mere fusion, a collapse of the many into an undifferentiated one, and differentiation without union would be a mere separation, a society of parts that never sound together. Sharing, on this model, is the musical moment, the bringing of the differentiated two into accord, and it presupposes rather than abolishes the differentiation that keeps them two; the concord is concord precisely because it is the accord of distinct voices and not the unison of one. This is the figure the section needs against the conflict ahead, and the tradition has held it for two and a half millennia: that the deepest union is the one in which the united remain distinct, that harmony is the sounding-together of difference and not its dissolution.

Returned to the light on the wall, the tradition reads the saying as the musical moment. Before the look, two persons differentiated, each in her own perception; the saying, look, and the looking that answers it, is the bringing of the two into accord upon the light, a concord in which they remain two who attend, the one who first saw and the one who is brought to see, joined in a focus that does not erase which of them first found it or how differently each receives it. The beauty produced is the accord itself, the sounding-together of two distinct perceivings upon one light, and it is real in the way a chord is real, present in the sounding-together and in neither note alone. That the tradition reached this without ever dividing the beautiful from the good is why it can state, as a single thought, what the Western accounts assemble from joint attention and witnessing and resonance: that to share the light is to tune, for a moment, the common life of the two, an act at once aesthetic and ethical because in this tradition those were never two acts.

\subsection{Two real conflicts}
\label{p22:sub:sharing-conflicts}

The accounts of sharing conflict at two points, the second of which is the dangerous one and the structural twin of the hardest steps already met.

The first conflict is the milder. Joint attention describes sharing as a coordination of attention, a recursive common focus; witnessing describes it as an ethical event, the lifting of a private occurrence into a held reality. Read as rivals, these would compete for what sharing fundamentally is, a matter of cognition or a matter of consequence. They are not rivals, and the framework declines to choose. The recursive common focus is the structure, and the witnessing is the weight that structure carries; a sharing can have the structure thinly, two attentions coordinated on an object with little at stake, or it can have the structure freighted with the consequence of accompaniment, and the difference between a thin coordination and a weighted witnessing is not a difference of two kinds of sharing but of how much the same structure bears. The framework therefore reads joint attention as the form of the between and witnessing as its ethical charge, and assigns the dispute to a confusion of levels rather than a genuine opposition. This is the easy conflict, recorded so that the hard one stands out.

The second conflict cuts to the concept. The accounts of joint attention, of dialogue, of musical concord all press toward union, the formation of a we, a common focus, a resonance in which two are mutually moved. Yet the reception section, following the welcome of the other, insisted that the other must be received as other and never reduced to the same, and a tradition descending from that welcome will warn that the we is the great solvent of otherness, that shared attention shades into fusion, dialogue into echo, resonance into a chamber in which each hears only a confirmation of self, and that the union sharing seeks may be the very dissolution of the difference on which relation depends. If sharing aims at union, and union dissolves the two, then sharing at its limit destroys what relational aesthetics requires, a beauty produced between two who remain two.

\subsection{The hardest step: sharing is grounded in the preservation of difference}
\label{p22:sub:sharing-hardest}

The resolution requires distinguishing two unions that the word conceals, and it is the same shape as the distinctions that resolved the earlier conflicts. One union is fusion, the collapse of two into one, in which difference is dissolved and what remains is a single enlarged self mistaking its own echo for an other. The other is concord, the harmony named in the musical tradition as union-without-uniformity, in which two remain two and precisely thereby can answer, resonate, and witness one another. Fusion cannot witness, because there is no longer a second to witness the first; it cannot answer, because answering requires an other who is not already oneself; it cannot resonate, because resonance is the mutual movement of two and not the ringing of one. The very phenomena the accounts of sharing describe, the recursive focus that requires two attendings, the answer that requires an addressee, the resonance that requires two poles, are possible only between two who are not fused. Fusion would abolish the conditions of the sharing it was mistaken for.

\claim{Sharing is grounded in the preservation of difference, not in its dissolution. The union that sharing seeks is concord, in which two remain two and thereby can attend to one another's attending, answer one another, and witness one another; it is not fusion, in which two collapse into one and the conditions of the sharing are abolished. To be witnessed confirms that I and the other are two, for a witness is precisely a second who is not the first. Healthy sharing therefore confirms otherness rather than consuming it, and the deepest sharing is the moment at which the difference between the two is most fully held.}

An objection arises from the side of love's own self-understanding, because lovers have always spoken of fusion as the height of union, of becoming one, of the dissolution of the boundary between two into a single shared being, and a theory that makes the preservation of difference the condition of sharing seems to deny what lovers most prize, the merging they take to be love's summit. The objection has the weight of a long tradition of erotic and mystical speech behind it, and it cannot be met by simply preferring difference to union. It is met by distinguishing the experience of fusion from its reality. Lovers do experience, in certain moments, something they describe as fusion, a heightened sense of union in which the boundary seems to dissolve, and the experience is real and precious and not to be denied. But the experience of boundary-dissolution is not the actual collapse of two into one, and the account can honor the experience while denying the metaphysics the objection draws from it. What occurs in the prized moment is not that two become one, which would abolish the very relation the lovers are in and leave a single being with no one to love, but that two who remain two reach a concord so complete that the experience of it feels like dissolution, a union so deep that it is felt as oneness though it remains, in reality, the concord of two. The feeling of fusion is the phenomenology of a profound concord, not the metaphysics of an actual merging, and the account claims only that the reality is concord, however the experience describes itself.

This matters practically and not only theoretically, because the belief that love aims at actual fusion does damage that the distinction prevents. A lover who believes union means becoming one will experience the other's persistent difference, her irreducible otherness, as a failure of love, an obstacle to the fusion that love is supposed to achieve, and will press toward a merging that, were it achieved, would destroy the relation by destroying one of its terms. The belief in fusion as love's aim licenses the absorption of one by the other under the name of love's fulfilment, the swallowing of the other's difference as the overcoming of a barrier, and this is the possessive sharing the section warned against, dressed now in the language of love's highest aspiration. The distinction between the experience of fusion and the reality of concord guards against this, by showing that the prized moment is the depth of concord and not the achievement of merging, so that the other's persistent difference is not the obstacle to union but its condition, the very thing that makes the concord a concord of two and not the solitude of one. To love well, on this account, is to seek the depth of concord that feels like fusion while honoring the difference that fusion would destroy, and the objection from love's self-understanding, met this way, confirms rather than refutes the preservation of difference, since what love prizes is correctly understood as a concord so deep it feels like oneness, and a concord requires the two that fusion would abolish.

So the second conflict, like its twins in the perception and reception sections, inverts. The union sharing seeks does not threaten the otherness the reception section guarded; rightly understood as concord rather than fusion, it requires and confirms that otherness, since every phenomenon that makes sharing an event presupposes two who remain distinct. The seeking of common ground while the differences are kept, which this series has named as its method, is here shown to be not a procedural preference but the very structure of healthy sharing: the common ground is the shared field, and the kept differences are the two who must remain two for the field to be shared at all. The hardest step of this section is therefore the same counsel once more, to bring into harmony without lording over and to unite without absorbing, read into the conditions of the between. A sharing that fused would be a sharing that possessed, one self swallowing the other under the name of love; a sharing that holds the difference is a sharing that generates, two who remain two producing between them a beauty neither could hold alone.

\subsection{A layered synthesis}
\label{p22:sub:sharing-synthesis}

The accounts of sharing take their places without identification. Joint attention is taken at the level of structure: it gives the form of the between, the recursive common focus that distinguishes an event from a coincidence of privacies, and it reaches toward the implementation section as the candidate for an observable trace. Witnessing is taken at the level of ethics: it gives the weight the structure carries, the lifting of the private into a held reality, the consequence of accompaniment against abandonment. Dialogue and resonance are taken at the level of dynamics: they give the movement by which the shared increases through answering, the bridge along which sharing passes into reproduction, and the contrast of resonance with reification that the political economy will need. The musical tradition is taken as the principal voice that states the thesis directly and supplies the figure of concord-without-uniformity, joining the aesthetic to the ethical and political without a seam. None is written as equal to another, because the form of the between is not its ethical charge, and neither is its dynamism, and to assert their identity would claim a structure none of them has.

What the synthesis shows is that sharing is the step at which relational aesthetics passes from a thesis about beauty into a visible event, and at which the method of the whole series, common ground with kept difference, is revealed as the inner structure of that event rather than a stance taken toward it. The between is born when two who remain two bring a field into a focus each knows to be shared, witness in one another the appropriateness that occurs there, and answer one another so that what occurs increases. That is the birth of the relational, and with it beauty becomes, at last, the coupling event the concept named, no longer a perception held in one subject but a thing produced and alive between two.

\subsection{The reach and the defeat of each account}
\label{p22:sub:sharing-ledger}

The accounts of sharing differ in what aspect of the between they grasp, and a frank accounting of their reach and their failure is owed before the layered arrangement can claim authority.

The account of joint attention explains the structure of the between, the recursive common focus that distinguishes a shared event from a coincidence of two privacies, and it alone among the accounts connects sharing to something that could be observed, the inter-brain coupling that a measure spanning two nervous systems might detect. What defeats it, taken alone, is its silence about why the between matters. It gives the form of sharing and not its weight, describing the recursive structure without saying what is at stake in it between the two, and a structure described without its stakes is a cognitive curiosity rather than the consequential event sharing is. Its strength is the form and its defeat is the silence about consequence.

The account of witnessing explains exactly the consequence that joint attention omits, why the between matters, what is at stake in being accompanied rather than left alone with experience, and it raises sharing from a coordination to an ethical event. What defeats it, taken alone, is its thinness about structure and dynamics. It says that being witnessed changes the standing of what is witnessed but says little about the recursive structure that makes witnessing possible or the exchange through which the witnessed beauty grows, and an account of sharing cannot rest on the ethical weight alone, needing the structure that joint attention supplies and the dynamics that the dialogical account supplies. Its strength is the weight and its defeat is the need for the form and the movement it does not provide.

The dialogical and sociological accounts explain the dynamics that the other two omit, the movement by which the shared beauty increases through answering and the mechanism by which co-present attention amplifies into a common reality, and they supply the bridge to reproduction and the caution that the common reality forged in sharing is morally indifferent. What defeats them, taken alone, is their indifference to the justice of what is produced, the effervescence of the gathered being as available to the mob as to the festival, so that they describe the production of a vivid common reality without distinguishing the just from the unjust production, which is why they must be bounded by the orthogonality the paper insists on. Their strength is the dynamics and their defeat is the silence about justice that the orthogonality must supply.

The materialist resolution holds that sharing is a real event between two real subjects in a real field, possessing at once a structure, an ethical weight, and a dynamics, and that the production of the common reality is a real material process whose justice is a further real question not settled by its vividness. It refuses to crown the structure, the weight, or the dynamics as the whole, holding each to be an aspect of one real event whose truth is their articulation. It is materialist in taking the between to be a real relation producing a real common reality, not a construction of either subject nor an appearance of an absolute; it is dialectical in holding the structure, the weight, and the dynamics to be aspects whose articulated relation is the truth of sharing. Against the idealism that would make the we a moment of one subject's self-relation, it holds the real difference of the two that concord preserves; against a reductive account that would make sharing mere behavioural coordination, it holds the ethical weight of witnessing; against an account that would celebrate the common reality without judging it, it holds the orthogonal question of justice. Sharing is the real production of a common reality between two who remain two, structured and weighted and dynamic and answerable to justice, and reducible to none of these, which is the materialist dialectic applied to the birth of the relational.
